\documentclass[9pt]{beamer}

\usecolortheme[RGB={205,173,0}]{structure} 
\usefonttheme{serif} 

\usepackage{amssymb}

%Typesetting conventions
\def\K3{\mathrm K3}
\def\CY#1{\mathrm{CY}_{#1}}
%\def\SU2{\mathrm{SU}(2)}
\def\SU#1{SU({#1})}
%\def\U1{\mathrm U(1)}
\def\U#1{U({#1})}
\def\SO#1{SO({#1})}
\def\O#1{O({#1})}
%multiplet and dimension
\def\mult #1#2{\mathrm{#1}[\![ {#2}]\!]}
%normal ordering
\def\normal#1{\,\!:\! #1 \!:\! \,}
%doublebar
\def\double #1{#1{\hbox{\kern-2pt $#1$}}}
%underline
\def\un#1{\underline #1}


%%%%%%%%%%%%%%%%%%%%%%%%%%%%%%%%%%%%%%%%%%%%%%%%
%%%%%%%%%%%%%%%%  Typesetting conventions  %%%%%%%%%%%%%%%%%%%

%doublebar
\def\double #1{#1{\hbox{\kern-2pt $#1$}}}
%Bras and kets
\def\ket#1{\left| #1\right>}
\def\bra#1{\left< #1\right|}
\def\braket#1#2{\left< #1\right|\left.#2\right>}
%For twistors
\def\kra#1{\left[ #1\right|}
\def\bet#1{\left| #1\right]}

%%%%%%%%%%%%%%%%%%%%%%%%%%%%%%%%%%%%%%%%%%%%%%%%
%%%%%%%%%%%%  Jim Gates and Marc Grisaru MACROS   %%%%%%%%%%%%%%%%
\def\pp{{\mathchoice
            %{general format
               %[w] = length of horizontal bars
               %[t] = thickness of the lines
               %[h] = length of the vertical line
               %[s] = spacing around the symbol
              %
              %\kern [s] pt%
              %\raise 1pt
              %\vbox{\hrule width [w] pt height [t] pt depth0pt
              %      \kern -([h]/3) pt
              %      \hbox{\kern ([w]-[t])/2 pt
              %            \vrule width [t] pt height [h] pt depth0pt
              %            }
              %      \kern -([h]/3) pt
              %      \hrule width [w] pt height [t] pt depth0pt}%
              %      \kern [s] pt
          {%displaystyle
              \kern 1pt%
              \raise 1pt
              \vbox{\hrule width5pt height0.4pt depth0pt
                    \kern -2pt
                    \hbox{\kern 2.3pt
                          \vrule width0.4pt height6pt depth0pt
                          }
                    \kern -2pt
                    \hrule width5pt height0.4pt depth0pt}%
                    \kern 1pt
           }
            {%textstyle
              \kern 1pt%
              \raise 1pt
              \vbox{\hrule width4.3pt height0.4pt depth0pt
                    \kern -1.8pt
                    \hbox{\kern 1.95pt
                          \vrule width0.4pt height5.4pt depth0pt
                          }
                    \kern -1.8pt
                    \hrule width4.3pt height0.4pt depth0pt}%
                    \kern 1pt
            }
            {%scriptstyle
              \kern 0.5pt%
              \raise 1pt
              \vbox{\hrule width4.0pt height0.3pt depth0pt
                    \kern -1.9pt  %[e]=0.15pt
                    \hbox{\kern 1.85pt
                          \vrule width0.3pt height5.7pt depth0pt
                          }
                    \kern -1.9pt
                    \hrule width4.0pt height0.3pt depth0pt}%
                    \kern 0.5pt
            }
            {%scriptscriptstyle
              \kern 0.5pt%
              \raise 1pt
              \vbox{\hrule width3.6pt height0.3pt depth0pt
                    \kern -1.5pt
                    \hbox{\kern 1.65pt
                          \vrule width0.3pt height4.5pt depth0pt
                          }
                    \kern -1.5pt
                    \hrule width3.6pt height0.3pt depth0pt}%
                    \kern 0.5pt%}
            }
        }}

\def\mm{{\mathchoice
                      %{general format %[w] = length of bars
                                       %[t] = thickness of bars
                                       %[g] = gap between bars
                                       %[s] = space around symbol
   %[w], [t], [s], [h]=3([g]) are taken from corresponding definitions of \pp
   %
                      %       \kern [s] pt
               %\raise 1pt    \vbox{\hrule width [w] pt height [t] pt depth0pt
               %                   \kern [g] pt
               %                   \hrule width [w] pt height[t] depth0pt}
               %              \kern [s] pt}
                  %
                       {%displaystyle
                             \kern 1pt
               \raise 1pt    \vbox{\hrule width5pt height0.4pt depth0pt
                                  \kern 2pt
                                  \hrule width5pt height0.4pt depth0pt}
                             \kern 1pt}
                       {%textstyle
                            \kern 1pt
               \raise 1pt \vbox{\hrule width4.3pt height0.4pt depth0pt
                                  \kern 1.8pt
                                  \hrule width4.3pt height0.4pt depth0pt}
                             \kern 1pt}
                       {%scriptstyle
                            \kern 0.5pt
               \raise 1pt
                            \vbox{\hrule width4.0pt height0.3pt depth0pt
                                  \kern 1.9pt
                                  \hrule width4.0pt height0.3pt depth0pt}
                            \kern 1pt}
                       {%scriptscriptstyle
                           \kern 0.5pt
             \raise 1pt  \vbox{\hrule width3.6pt height0.3pt depth0pt
                                  \kern 1.5pt
                                  \hrule width3.6pt height0.3pt depth0pt}
                           \kern 0.5pt}
                       }}

\def\ad{{\kern0.5pt
                   \alpha \kern-5.05pt
\raise5.8pt\hbox{$\textstyle.$}\kern 0.5pt}}

\def\bd{{\kern0.5pt
                   \beta \kern-5.05pt \raise5.8pt\hbox{$\textstyle.$}\kern 0.5pt}}

\def\qd{{\kern0.5pt
                   q \kern-5.05pt \raise5.8pt\hbox{$\textstyle.$}\kern 0.5pt}}
\def\Dot#1{{\kern0.5pt
     {#1} \kern-5.05pt \raise5.8pt\hbox{$\textstyle.$}\kern 0.5pt}}
%%%%%%%%%%%%%%%%%%%%%%%%%%%%%%%%%%%%%%%%%%%%%%%%%%%
%%%%%%%%%%%%%%%%%%%%%%%%%%%%%%%%%%%%%%%%%%%%%%%%%%%


\usepackage{beamerthemesplit}

\title{The Covariant Superstring on K3}
\author{William D. Linch, III}
\date{Prepared for the Max Ba\~nados Field Theory and Gravitational Group,\\
Pontificia Universidad Cat\'olica de Chile,\\
3 November, 2011.
}

\begin{document}

\frame{\titlepage}

\section[Outline]{}
\frame{\tableofcontents}

%%%%%%%%%%%%%%%%%%%%%%%%%%%%%%%%%%%%%%%%%%%%%%%%%%%
\section{Introduction}
\frame{
\frametitle{Introduction}
\begin{itemize}
\item We want to use superstrings as a calculus for quantum field theories coupled to quantum gravity.
\item The simplest definition of this calculus is as a free ten-dimensional string.
\item To use this tool for lower-dimensional physics, we need to reduce from 10D.
\item First method: Holographic (Brenno's talk here last week)
\item Second method: Compactification
\item The compactification problem of relating the pure spinor in 10D to the RNS string in lower dimensions has resisted solution since 2002.
\item We will solve this problem for compactification of the pure spinor string to six dimensions on a K3 surface ({\it i.e.} the four-dimensional Calabi-Yau manifold).
\end{itemize}
}
%%%%%%%%%%%%%%%%%%%%%%%%%%%%%%%%%%%%%%%%%%%%%%%%%%%
\section{Covariant superstrings}
\frame{
\frametitle{The Green-Schwarz Superstring (1984)}
We start with the bosonic string $x:\Sigma \to M$ with dynamics determined by
\begin{eqnarray}
S_\mathrm {bose}\sim \int_\Sigma (\partial x)^2.
\nonumber
\end{eqnarray}
\pause
Then we add space-time ($M$) supersymmetry
\begin{eqnarray}
\delta x^m = i\epsilon \gamma^m \theta,~~~
\delta \theta=\epsilon.
\nonumber
\end{eqnarray}
\pause
The Green-Schwarz variables are manifestly supersymmetric
\begin{eqnarray}
\Pi^m = \mathrm dx^m-i \theta \gamma^m \mathrm d \theta,~~~
\mathrm d\theta.
\nonumber
\end{eqnarray}
\pause
And the Green-Schwarz superstring is defined by the action
\begin{eqnarray}
S_\mathrm{GS} \sim \int_\Sigma \Pi^2 + S_\mathrm{WZ}.&&\cr
\pause
\uparrow~~&&\cr
\textrm{(for Max)}&&
\nonumber
\end{eqnarray}
\pause
%Additionally, there is a Siegel symmetry
%\begin{eqnarray}
%\delta \Pi= \mathrm d\theta\gamma^m \delta \theta,~~~
%\delta \theta = \slash\!\!\! \Pi \kappa.
%\nonumber
%\end{eqnarray}
Unfortunately the flat space limit of this string is only free in light-cone gauge so covariant quantization is not straightforward.
}

%%%%%%%%%%%%%%%%%%%%%%%%%%%%%%%%%%%%%%%%%%%%%%%%%%%
\frame{
\frametitle{Siegel formalism (1986)}
Siegel proposed the following modification:
\begin{eqnarray}
S_\mathrm{GSS} \sim S_\mathrm{GS}+ \int_\Sigma d\partial \theta,
\nonumber
\end{eqnarray}
where the Siegel constraint is given by
\begin{eqnarray}
d= p_\theta + i (\theta \gamma^m)\partial x_m + \theta^2\partial \theta\textrm{-term}.
\nonumber
\end{eqnarray}
\pause
With this, the string in flat space is free
\begin{eqnarray}
S_\mathrm{GSS} \sim \int_\Sigma \left[ (\partial x)^2+ p_\theta \partial \theta\right],
\nonumber
\end{eqnarray}
\pause
So quantization gives the canonical operator product algebra
\begin{eqnarray}
x(z) x(0) \sim \log|z|^2, ~~~ p_\theta(z) \theta(0)\sim \frac 1z.
\nonumber
\end{eqnarray}
}

%%%%%%%%%%%%%%%%%%%%%%%%%%%%%%%%%%%%%%%%%%%%%%%%%%%
\frame{
\frametitle{Dirac quantization}
Then it is easy to compute the operator product
%The fundamental operator product relation of the Siegel constraints is
\begin{eqnarray}
d(z)d(0) \sim \frac1z \slash\!\!\!\Pi(0),
\nonumber
\end{eqnarray}
%\pause
implying that (half of) the Siegel constraints are $2^{\textrm{nd}}$-class constraints.\\
\pause
There is also the $1^{\textrm{st}}$-class Virasoro constraint
\begin{eqnarray}
T \sim\, \normal{\Pi \Pi}+\normal{d\partial \theta}
\nonumber
\end{eqnarray}
\pause
Problem: To use the Dirac procedure, we have to separate the constraints. Is there a Poincar\'e-covariant way to do this?
}

%%%%%%%%%%%%%%%%%%%%%%%%%%%%%%%%%%%%%%%%%%%%%%%%%%%
\frame{
\frametitle{Berkovits' formalism I (1994): Hybrid}
Consider the RNS string on a factorized background
\begin{eqnarray}
\begin{array}{ccc}
M&\to& M\times X\\
\mathrm{Hol}(\nabla_M)&\to&\mathrm{Hol}(\nabla_M)\times\mathrm{Hol}(\nabla_X)
\end{array}
\nonumber
\end{eqnarray}
\pause
If $\mathrm{Hol}(\nabla_X)$ is ``special'' 
\pause
then $\exists$ map
\begin{eqnarray}
\begin{array}{lcc}
	&\textrm{``space-time''}& \textrm{``internal''}\\\\
\textrm{matter}&x^m , \psi^m & y^i , \bar y^{\bar \imath}, \psi^i, \bar \psi^{\bar \imath}\\
\textrm{ghosts} & b,c,\beta,\gamma &\\
&\downarrow&\downarrow\\
\textrm{Berkovits}~
\textrm{variables}&x^m ,p_\theta, \theta,\rho,\sigma & y^i , \bar y^{\bar \imath}, \psi^i, \bar \psi^{\bar \imath}\\
\end{array}
\nonumber
\end{eqnarray}
with canonical operator products.\pause..\\
and a factorized algebra of constraints
\begin{eqnarray}
\mathrm {SCA}_{10} \to \mathrm {SCA}_{\textrm{space-time}}\oplus \mathrm {SCA}_{\textrm{internal}}.
\nonumber
\end{eqnarray}
}

%%%%%%%%%%%%%%%%%%%%%%%%%%%%%%%%%%%%%%%%%%%%%%%%%%%
\frame{
\frametitle{Berkovits' formalism II (2000): Pure spinor}
The maximal number of supercharges realizable in this way is 8.\\
\pause
Meanwhile, the minimal number of supersymmetries in $D=10$ is 16.\\~\\
\pause 
Therefore, a new method is required to separate the ten-dimensional constraints
\begin{eqnarray}
%d_\alpha (z) d_\beta(0)\sim \frac 1z (\gamma^a)_{\alpha \beta} \Pi_m(0)
d(z)d(0) \sim \frac1z \slash\!\!\!\Pi(0).
\nonumber
\end{eqnarray}
\pause
Borrow trick from harmonic superspace:
\pause
Introduce bosonic ``pure'' spinor $\lambda$:
\begin{eqnarray}
(\lambda \gamma^m \lambda)=0 ~~~ m=0,1,\dots,9.
\nonumber
\end{eqnarray}
\pause
(Together with their momenta, the pure spinors form a 22-dimensional cone over the space $SO(10)/U(5)$ of orthogonal complex structures on $\mathbb R^{10}$.)\\
\pause
Then define Berkovits differential 
\begin{eqnarray}
Q=\oint \lambda d ~~~:~~~ Q^2=0.
\nonumber
\end{eqnarray}
\pause
Proceed with BRST-type quantization!
}

%%%%%%%%%%%%%%%%%%%%%%%%%%%%%%%%%%%%%%%%%%%%%%%%%%%
\frame{
\frametitle{Mysteries}
\begin{itemize}
\item What is the pure spinor?
%\pause
Harmonic? Ghost? Twistor? Mix?\\
%\pause
%Looks like gauge-fixed superstring...\\
%\pause
\item What is the relation with the hybrid formalism?\\
%\pause
\item What is the relation (if any) with the pure spinors of generalized geometry?\\
%\pause
\item How to recover the correct spectrum in compactification with off-shell superspace? \\
\item Many more questions...
\end{itemize}
\pause
Strategy: Compactify on K3 
\pause
which has the advantages that it
\begin{itemize}
\item is simple,\pause
\item has a hybrid description, \pause
\item has harmonics in its superspace, \pause
\item can be compactified further to 4 dimensions on $T^2$ to make K3-fibered Calabi-Yau 3-fold.
\end{itemize}
}


%%%%%%%%%%%%%%%%%%%%%%%%%%%%%%%%%%%%%%%%%%%%%%%%%%%
\section{Compactification}
\frame{
\frametitle{K3 orbifold}
Start with a torus ({\em very} special holonomy)
\begin{eqnarray}
T^4&\approx& {\mathbb C \times \mathbb C\over \mathbb Z\times \mathbb Z},
%\nonumber
%\end{eqnarray}
%\pause
%Hodge diamond 
%\begin{eqnarray}
~~~~~~~~~
\begin{array}{ccccc}
&&1&&\\
&2&&2&\\
1&&4&&1\\
&2&&2&\\
&&1&&
\end{array}
\nonumber
\end{eqnarray}
\pause
and gauge $\mathbb Z_2\subsetneq SU(2)_- \hookrightarrow SU(2)_+\times SU(2)_- \approx Spin(4)$ subgroup of the four-dimensional Lorentz ({\it i.e.} rotation) group.\\
\pause
%Geometrically, this identifies $(z^1,z^2) \sim (-z^1,-z^2)$ to get the K3 orbifold 
\begin{eqnarray}
\Rightarrow~~~ 
\textrm{K3 orbifold } X=T^4/\sim,
~~~~~~~~~
\begin{array}{ccccc}
&&1&&\\
&0&&0&\\
1&&4&&1\\
&0&&0&\\
&&1&&
\end{array}
\nonumber
\end{eqnarray}
with $2^4=16$ fixed points
\pause
which can be blown up to Eguchi-Hanson spaces $\mathcal O_{\mathbb P^1}(-2)\approx T^*S^2$ (gravitational instantons).
}

\frame{
\frametitle{Twisted sector states}
Orbifold identification gives new massless ``twisted'' string states at each orbifold point.\\
\pause
The Casimir energies of the worldsheet fields is given by
\begin{eqnarray}
E= (-)^F\left[ \frac1{48} -\frac1{16}(2\theta-1)^2\right]
\left\{ 
\begin{array}{rl}
-\frac1{24}&{\textrm {real periodic boson}}\\~\\
\frac1{48}&{\textrm {real anti-periodic boson}}
\end{array}
\right.
\nonumber
\end{eqnarray}
\pause
Contribution with $n$ untwisted pure spinors:
\pause
\begin{eqnarray}
E&=&\underbrace{6( -\tfrac1{24})+ 4(\tfrac1{48})}_{x+\tilde x}+
\underbrace{8(\tfrac1{12})+ 8(-\tfrac1{24})}_{\theta, p+ \tilde \theta ,\tilde p}+
\underbrace{n(-\tfrac1{12})+(11-n)(\tfrac1{24})}_{\lambda,w+ \tilde \lambda, \tilde w}\cr
&=&-\frac 18 (n-5).
\nonumber
\end{eqnarray}
\pause
Spectrum matches iff $n=5\Rightarrow6$ of the original 11 pure spinors are $\mathbb Z_2$-twisted.
}

%%%%%%%%%%%%%%%%%%%%%%%%%%%%%%%%%%%%%%%%%%%%%%%%%%%
\section{Dimensional reduction of the pure spinor}
\frame{
\frametitle{$U(5)$ variables}
Breaking $Spin(10)\to U(5)$, we can solve the pure spinor constraint
\begin{eqnarray}
\begin{array}{ccccccc}
\lambda & \to & \gamma & &\gamma u_{[ab]} && \gamma u_{[ab}u_{cd]}\cr\cr
\mathbf{16}&=&\mathbf{1}&\oplus&\overline{\mathbf{10}}&\oplus&\mathbf{5}
\end{array}
\nonumber
\end{eqnarray}
\pause
Dimensionally reducing further $SU(5)\to SU(3)\times SU(2)$, we find that 
\begin{eqnarray}
\begin{array}{ccccccc}
\overline{\mathbf{10}} &=&(\mathbf{3},\mathbf 1)&\oplus&(\bar{\mathbf{3}},\mathbf2)&\oplus&(\mathbf 1,\mathbf{1})\\\\
u_{[ab]} &=&u^I&,&u_{Ia^\prime}&,&u.
\end{array}
\nonumber
\end{eqnarray}
\pause
This can be rearranged into a six-dimensional spinor $(\lambda^A)=\left(\begin{array}{c}\gamma\\\gamma u^I\end{array}\right)$, 
\pause
a ``harmonic'' iso-spinor: $(u^a)=\left(\begin{array}{c}1\\ u\end{array}\right)$, and
\pause
junk: $(u_{Ia^\prime}) \leftrightarrow \mathbf6$.
}

\frame{
\frametitle{Small Hilbert Space}
The Berkovits differential splits
\begin{eqnarray}
Q=\oint \left( \lambda^Au^a d_{Aa} +\lambda_{Aa^\prime} d^{Aa^\prime} + \Lambda^A d_{A2}\right)=Q_0+Q_1+Q_2
\nonumber
\end{eqnarray}
were $Q_m$ has $m$ factors of $\mathbf 6$ in it.\\
\pause
The Hilbert space of string fields forms a chain complex with differential $Q$ graded by number of $\mathbf 6$s
\begin{eqnarray}
\cdots \stackrel Q\longrightarrow C^\bullet \stackrel Q\longrightarrow C^\bullet \stackrel Q\longrightarrow 
C^\bullet \stackrel Q\longrightarrow \cdots
\nonumber
\end{eqnarray}
\pause
Claim: 
\begin{eqnarray}
H_Q (C^\bullet) = H^0_{Q_0}(C^\bullet).
\nonumber
\end{eqnarray}
In words: The physical states are described by the cohomology of a short differential on a small Hilbert space without $\mathbf 6$-variables.
}

\frame{
\frametitle{Spectrum}
The reduced differential factorizes the cohomology further:
\begin{eqnarray}
Q_0 = q+\delta
\nonumber
\end{eqnarray}
\pause
were 
\begin{eqnarray}
q=\oint \lambda^Ad^+_A~~~\mathrm{and}~~~ \delta = \oint \psi_{aa^\prime}\partial x^{aa^\prime} 
\nonumber
\end{eqnarray}
are both differentials.
Here,
%\begin{eqnarray}
$d^+_{A} = u^a\left(p_{Aa} + i\theta^B_a \partial x_{AB}+\epsilon_{ABCD} \theta^{B}_{a}\theta^{Cb}\partial \theta^D_b\right)$
%\nonumber
%\end{eqnarray}
and 
%\begin{eqnarray}
$\psi_{aa^\prime}= \gamma \,u_{a}\, \theta_{1a^\prime}$.\\
%\nonumber
%\end{eqnarray}
\pause
By Dolbeault's theorem,
%a standard argument in homological algebra, 
the cohomology of the double complex becomes the relative cohomology
\begin{eqnarray}
H_{Q_0} \approx H_{\delta}(H_{q}).
\nonumber
\end{eqnarray}
}

\frame{
\frametitle{Spectrum}
We have found that the spectrum of the pure spinor formalism on K3 is computed by the relative cohomology $H_{\delta}(H_{q})$.\\~\\
\pause
We will review below that $H_q$ is the space of projective superfields. \\~\\
\pause
Meanwhile, $H_\delta$ is the ring of vertex operators of the form $\omega_{aa^\prime bb^\prime\dots}\psi^{aa^\prime}\psi^{bb^\prime}\dots$ which are annihilated by $\delta= \psi^{aa^\prime}\partial_{aa^\prime}$ and not $\delta$-exact.\\~\\
\pause
In other words, $H_\delta \approx H_\mathrm{dR}$ is just the de Rham cohomology of K3.\\~\\
\pause
Hence, $H_\delta(H_q)$ is the cohomology algebra of K3 vertex operators with projective superfield wave functions.
}


\frame{
\frametitle{Projective superspace}
The space-time part of the hybrid formalism on K3 lives in projective superspace 
\begin{eqnarray}
\mathbb R^{6|8}\times \mathbb CP^1 &\approx& {\mathbb R^{6|8}\rtimes \left(Spin(6)\times SU_R(2)\right)\over
	Spin(6) \times U(1)_R}.
\nonumber
\end{eqnarray}
\pause
Vertex operators have superspace wave functions $\Phi(x,\theta,u)$,
%\pause
which can be expanded in harmonics
\begin{eqnarray}
\Phi(x,\theta,u) = \sum_{n=-\infty}^\infty u^n \varphi_n(x,\theta) 
\nonumber
\end{eqnarray}
%\pause
(infinite number of ordinary $\mathcal N=(1,0)$ superfields).\\~\\
\pause
The Siegel constraints act, by the operator product, as superspace derivatives
\begin{eqnarray}
d_{\alpha i} (z) \Phi(x,\theta, u) =\frac1z D_{\alpha i} \Phi(x,\theta,u) .
\nonumber
\end{eqnarray}
}

\frame{
\frametitle{Projective superspace}
The flat, six-dimensional superspace derivatives satisfy
\begin{eqnarray}
\{ D_{\alpha i} , D_{\beta j} \} = 2i \varepsilon_{ij} (\gamma^a)_{\alpha \beta} \partial_a.
\nonumber
\end{eqnarray}
\pause
Introduce zwei-beine $\{u^{\pm}_i\}_{i=1}^2$ of $S^2$ and define $D^\pm_{\alpha} = u^{\pm i} D_{\alpha i}$. 
\pause
Then
\begin{eqnarray}
\{ D^+ , D^+ \} = 0
\nonumber
\end{eqnarray}
gives the integrability condition for the ``analytic subspace'' $\mathrm {ker} D^+$.\\
\pause
The space parameterized by $(x,\theta^\pm, u^\pm)$ is called ``harmonic superspace''.\\
\pause
The analytic subspace parameterized by $(x, \theta^+,u^+)$ and in $\mathrm {ker} D^+$ is called ``projective superspace''.\\
\pause
Returning to the cohomology, the differential
\begin{eqnarray}
q=\oint \lambda d^+
\nonumber
\end{eqnarray}
\pause
acts as $D^+$, implying that string wave functions are projective superfields.
}

\frame{
\frametitle{Solving the off-shell problem}
Projective superspace is off-shell: Equations of motion not imposed by supersymmetry. \\~\\
\pause
In fact, $H_q=\{0\}$ if the superfield $\Phi$ is not entire on $\mathbb CP^1$.\\~\\
\pause
Restriction of $\Phi$ to be an entire function truncates the harmonic expansion and this puts the component superfields on-shell.\\~\\
\pause
To impose this in string theory, we need to bosonize the projective parameter:	
\pause
\begin{eqnarray}
u = \mathrm e^{-\rho -i\sigma}
\nonumber
\end{eqnarray}
in terms of chiral bosons
\begin{eqnarray}
\rho(z)\rho(0) \sim -\log z,~~~~~ \sigma(z)\sigma(0)\sim -\log z.
\nonumber
\end{eqnarray}
\pause
To ensure only $\rho/\sigma$-dependence in the combination $u$, vertex operators must have no poles with the constraint $J= \partial(\rho+i\sigma)$ (which generates shifts $(\rho, \sigma)\mapsto (\rho+a, \sigma+ia)$).
}

\frame{
\frametitle{Emergence of worldsheet superconformal symmetry}
The condition that the harmonic expansion of the vertex has no poles is implied by the constraint
\begin{eqnarray}
G^- = \normal{\mathrm e^{-i\sigma}}.
\nonumber
\end{eqnarray}
\pause
Together with the Virasoro constraint 
\begin{eqnarray}
T=\normal{ \Pi^2}+ \normal{d^+\partial \theta^-}+\dots,
\nonumber
\end{eqnarray}
\pause
the algebra closes if we also define
\begin{eqnarray}
G^+ = -\frac1{4!} \epsilon^{ABCD}\normal{ d^+_A(d^+_B(d^+_C(d^+_D(\mathrm e^{2\rho+3i\sigma}))))}.
\nonumber
\end{eqnarray}
\pause
This terminates the expansion and puts the wave function on-shell.\\
\pause
The constraints $\{J, G^\pm, T\}$ generate the desired $N=2$ worldsheet superconformal algebra of the {\em hybrid} (a.k.a. RNS) superstring on K3.
}


%%%%%%%%%%%%%%%%%%%%%%%%%%%%%%%%%%%%%%%%%%%%%%%%%%%
\section{Conclusion}
\frame{
\frametitle{R\'esum\'e} 
\begin{itemize}
\item We started with the pure spinor superstring in ten flat dimensions.
\pause
\item We solved the pure spinor condition and rewrote all worldsheet fields in $(6+2+2^*)$-variables appropriate to compactification on K3.
\pause
\item We proved a theorem that the cohomology reduces to a small Hilbert space which removes 6 of the original 11 pure spinor degrees of freedom.
\pause
\item We interpreted the remaining $4+1$ pure spinor variables as an unconstrained $Spin(5,1)$ spinor and a holomorphic projective variable $u$ parameterizing $\mathbb CP^1\approx SU(2)/U(1)$.
\pause
\item We bosonized $u$ and defined the constraints $J$ and $G^-$ which simplify the description of the physical state conditions.
\pause
\item We closed the constraint algebra by introducing $G^+$, thereby recovering the physical state conditions.
\pause 
\item Thus, we succeed in deriving the hybrid string on K3 from the compactification of the ten-dimensional pure spinor string.\pause Yay!
\end{itemize}
}


%%%%%%%%%%%%%%%%%%%%%%%%%%%%%%%%%%%%%%%%%%%%%%%%%%%

\frame{
\frametitle{Fin}
\begin{centering}
{\Huge Thank you!}
\end{centering}
}

%%%%%%%%%%%%%%%%%%%%%%%%%%%%%%%%%%%%%%%%%%%%%%%
The work reported here was done in collaboration with Osvaldo Chand\'ia and Brenno Carlini Vallilo in the following papers: 
\begin{thebibliography}{99}
\footnotesize{

%%\bibitem{AP13}
%%\cite{Linch:2006ig}
%\bibitem{Linch:2006ig}
%  Wm~D.~Linch, III and B.~C.~Vallilo,
%  ``Hybrid formalism, supersymmetry reduction, and Ramond-Ramond fluxes,''
%  JHEP {\bf 0701}, 099 (2007)
%  [arXiv:hep-th/0607122].
%  %%CITATION = JHEPA,0701,099;%%
  
%%\cite{Berkovits:1999du}
%\bibitem{Berkovits:1999du}
%  N.~Berkovits,
%  ``Quantization of the type II superstring in a curved six-dimensional
%  background,''
%  Nucl.\ Phys.\  B {\bf 565}, 333 (2000)
%  [arXiv:hep-th/9908041].
%  %%CITATION = NUPHA,B565,333;%%

%%\cite{Linch:2006sh}
%\bibitem{Linch:2006sh}
%  Wm~D.~Linch, III and B.~C.~Vallilo,
%  ``Covariant N = 2 heterotic string in four dimensions,''
%  JHEP {\bf 0703}, 082 (2007)
%  [arXiv:hep-th/0611105].
%  %%CITATION = JHEPA,0703,082;%%
  
%%\cite{Linch:2008rw}
%\bibitem{Linch:2008rw}
%  Wm~D.~Linch, III, J.~McOrist and B.~C.~Vallilo,
%  ``Type IIB Flux Vacua from the String Worldsheet,''
%  JHEP {\bf 0809}, 042 (2008)
%  [arXiv:0804.0613 [hep-th]].
%  %%CITATION = JHEPA,0809,042;%% 

%%\cite{Linch:2008nt}
%\bibitem{Linch:2008nt}
%  Wm~D.~Linch, III and B.~C.~Vallilo,
%  ``Integrability of the Gauged Linear Sigma Model for $AdS_5\times S^5$''
%  JHEP {\bf 0911}, 007 (2009)
%  [arXiv:0804.4507 [hep-th]].
%  %%CITATION = JHEPA,0911,007;%%
  
%\cite{Chandia:2009it}
\bibitem{Chandia:2009it}
  O.~Chandia, Wm~D.~Linch III and B.~C.~Vallilo,
  ``Compactification of the Heterotic Pure Spinor Superstring I,''
  JHEP {\bf 0910}, 060 (2009)
  [arXiv:0907.2247 [hep-th]].
  %%CITATION = JHEPA,0910,060;%%
  
%\cite{Chandia:2011wd}
\bibitem{Chandia:2011wd}
  O.~Chandia, Wm~D.~Linch III, B.~C.~Vallilo,
  ``Compactification of the Heterotic Pure Spinor Superstring II,''
  JHEP {\bf 1110}, 098 (2011).
  [arXiv:1108.3555 [hep-th]].
  
%\cite{Chandia:2011su}
\bibitem{Chandia:2011su}
  O.~Chandia, Wm~D.~Linch III, B.~C.~Vallilo,
  ``The Covariant Superstring on K3,''
  [arXiv:1109.3200 [hep-th]].

}%end footnotesize

\end{thebibliography}


\end{document}